\documentclass[12pt]{article}
\usepackage[utf8]{inputenc}
\usepackage[russian]{babel}
\usepackage{amsmath, amssymb}
\usepackage{graphicx}
\usepackage{geometry}
\usepackage{indentfirst}
\usepackage{cite}
\usepackage{hyperref}
\usepackage{csvsimple}
\usepackage{booktabs}
\usepackage{caption}
\usepackage{etoolbox}
\usepackage{float}
\geometry{a4paper, margin=2.5cm}

\begin{document}

\pagenumbering{gobble}
\begin{titlepage}
    \centering
    \textbf{Московский государственный университет имени М.В. Ломоносова\\
    Филиал в городе Ташкенте\\
    Механико-математический факультет\\
    Кафедра Математической теории интеллектуальных систем\\}
    \vfill
    \textbf{\Large{Выпускная квалификационная работа}}\\
    \vspace{0.5cm}
    по направлению «Прикладная математика и информатика»\\
    на тему: \\
    \vspace{0.5cm}
    \textbf{\Large{Автоматная модель психологического взаимодействия}}\\
    \vfill
    \begin{flushright}
        \textbf{Подготовил:} студент четвертого курса ПМиИ\\ Таджиев Георгий Тимурович\\
        \textbf{Научный руководитель:} старший научный сотрудник\\ Волков Николай Юрьевич\\
    \end{flushright}
    \vfill
    Ташкент — 2026
\end{titlepage}
\pagenumbering{arabic}
\clearpage

\begin{abstract}
В данной работе рассматривается дискретная автоматная модель психологического взаимодействия агентов в макро-социальных группах (коллективах масштаба университета). В основе модели лежит формализация эмоциональных состояний и реляционных связей, позволяющая исследовать динамику социальных структур. Произведена масштабная архитектурная оптимизация гибридного вычислительного ядра (Python/C++), позволившая осуществлять симуляции коллективов, состоящих из более чем полутора тысяч агентов. Для обеспечения аналитики класса Big Data внедрена колоночная СУБД ClickHouse. Описана методология проведения многосерийных вычислительных экспериментов.
\end{abstract}

\tableofcontents
\clearpage

\section{Введение}

Моделирование сложных социальных систем с помощью многоагентных симуляций становится важным инструментом для изучения динамики взаимодействий между индивидами. В частности, учет эмоциональных состояний и реляционных функций позволяет более точно описывать процессы принятия решений и формирования отношений в обществе. Переход от изучения микро-групп к моделированию макро-социальных структур (например, социальной экосистемы целого университета на тысячи агентов) требует как концептуального, так и глубокого архитектурного переосмысления существующих подходов.

Современные исследования в области межагентных систем требуют междисциплинарного подхода, где психология, информатика и социология взаимодействуют для создания более реалистичных моделей. Психологические теории помогают понять эмоциональные и когнитивные механизмы агентов, информатика обеспечивает алгоритмическую реализацию и обработку огромных массивов данных, а социология предлагает концептуальные рамки для анализа социальных структур и динамики. 

Особую значимость в этой области приобрели исследования по аффективным вычислениям. Как отмечала Розалинд Пикар и ее последователи в исследованиях по аффективным вычислениям, учет эмоций при проектировании интеллектуальных систем является неотъемлемой частью создания правдоподобного поведения искусственных агентов. Аналогично, работы по когнитивной социологии, в частности, когнитивная модель доверия Кастельфранчи и Фальконе, подтвердили, что доверие и эмоциональный интеллект являются базисом для формирования стабильных межличностных графов. Согласно этой модели, доверие рассматривается как сложный когнитивный процесс, включающий оценку надежности и доброжелательности другого субъекта на основе имеющейся информации и прошлых взаимодействий.

Целью данной работы является разработка и программная реализация высокопроизводительной автоматной модели психологического взаимодействия. Это предоставляет инструмент для изучения комплексных социальных процессов, где эмоции и отношения взаимно влияют друг на друга, и позволяет исследовать механизмы устойчивости и изменения социальных связей на масштабах, приближенных к реальным коллективам.

Теоретической основой данной модели являются представления из теории автоматов, в частности концепции автоматных коллективов. Агенты в рассматриваемой модели интерпретируются как дискретные автоматы, переходы которых зависят от входных состояний и реляционных параметров. Такой подход согласуется с идеями, изложенными в трудах В. Б. Кудрявцева и соавторов [1], а также с моделью автоматного взаимодействия, представленной Н. Ю. Волковым [2].

\section{Определения}

В данном разделе собраны ключевые формальные и семантические определения, используемые в модели.

\paragraph{Агент (\(\mathcal{A}_i\))} — базовый дискретный автомат, обладающий внутренним состоянием (эмоциями), правилами перехода и набором направленных связей к другим агентам.
\paragraph{Коллектив (\(\mathcal{K}\))} — множество агентов \( \mathcal{K} = \{\mathcal{A}_1, \mathcal{A}_2, \ldots, \mathcal{A}_n \} \), находящихся в общем социальном пространстве и способных вступать во взаимодействия друг с другом.

\paragraph{Эмоциональные оси} — независимые параметры, описывающие психологический фон агента. Мы исходим из предположения, что каждая ось эмоций является статистически независимой. Оси включают:
\begin{itemize}
    \item \textbf{Радость — Печаль}: уровень общего удовлетворения или подавленности.
    \item \textbf{Страх — Спокойствие}: уровень тревожности и ожидания угрозы.
    \item \textbf{Гнев — Смирение}: агрессивность и склонность к конфликтному поведению.
    \item \textbf{Отвращение — Принятие}: реакция на социальные и внутренние стимулы.
    \item \textbf{Удивление — Привычка}: реакция на новизну и нарушение рутины.
    \item \textbf{Стыд — Уверенность}: уровень социальной самооценки агента.
    \item \textbf{Открытость — Отчуждение}: готовность к социальной интеграции.
\end{itemize}

\paragraph{Реляционные функции} — метрики, отражающие направленное отношение агента-субъекта к агенту-объекту:
\begin{itemize}
    \item \textbf{Выгода (\(U_{ij}\))} — прагматическая полезность взаимодействия. Отражает степень результата, который получает агент с точки зрения личных интересов.
    \item \textbf{Симпатия (\(A_{ij}\))} — иррациональное или эмоциональное расположение, не обязательно связанное с рациональной оценкой поведения партнёра.
    \item \textbf{Доверие (\(T_{ij}\))} — когнитивная оценка предсказуемости. Это ожидание, что партнёр будет вести себя предсказуемо и доброжелательно, обеспечивая стабильность взаимодействия.
\end{itemize}

\paragraph{Чувствительность (\(s_i\))} — скалярный коэффициент в интервале \((0, 3]\), определяющий восприимчивость к изменениям.

\paragraph{Контекст среды (\(\epsilon\))} — динамический коэффициент местонахождения, отражающий, как конкретная локация (аудитория, спортзал, коридор) влияет на социальное поведение агента.

\paragraph{Параметры посещаемости} — индивидуальные маркеры поведения агента:
\begin{itemize}
    \item \textbf{Склонность к пропускам (\(\chi_i\))} — случайная величина, определяющая вероятность агента пропустить академические занятия или покинуть университет досрочно.
    \item \textbf{Спортивность (\(\rho_i\))} — случайная величина, определяющая вероятность агента посетить спортивный зал или одну из трех добровольных спортивных секций.
\end{itemize}

\paragraph{Исход взаимодействия (\(\sigma_{i,j}\))} — дискретная переменная акта коммуникации, которая определяется как:
\[
\sigma_{i,j} = \begin{cases} 
1, & \text{успешный контакт,} \\
0, & \text{отказ от взаимодействия,} \\
-1, & \text{неудачный контакт.}
\end{cases}
\]
При состоявшемся контакте (когда не произошло целеустремленного отказа), исходы \(1\) и \(-1\) реализуются с равной вероятностью (\(p=0.5\)) и служат для моделирования влияния внешних стохастических факторов, не учтенных в детерминированной части автомата. При исходе \(\sigma_{i,j} = 1\) происходит положительное скалярное смещение: позитивные реляционные связи укрепляются, а негативные ослабевают. При исходе \(\sigma_{i,j} = -1\) происходит обратный процесс: позитивные связи ослабевают, а негативные усугубляются. В случае же отказа (\(\sigma_{i,j} = 0\)) оба участника корректируют реляционные функции друг к другу строго отрицательным образом в соответствии со своими архетипами.

\paragraph{Архетипы агентов}
Архетип — это фиксированная стратегия поведения автомата. Он задает нелинейные трансформации расчета приоритета, базовый шанс отказа, температуру принятия решений (\(\tau\)) и адаптивность к среде (\(\epsilon\)). Модель приоритета строится как линейная сумма нелинейных преобразований реляционных функций: \(P_{ij} = \epsilon \cdot (f_A(A) + f_U(U) + f_T(T))\).

\begin{itemize}
    \item \textbf{Аналитик} — хладнокровный, избегает крайностей, эмоционально стабилен, ценит знание и логику. Реже инициирует эмоциональные реакции, проявляет рациональность в выборе взаимодействий. 
    \textit{Особенности}: Формула приоритета \(P_{ij} = \epsilon \cdot (A_{ij} + \log(U_{ij}) + T_{ij})\). Адаптивность к среде максимальна на учебе (\(\epsilon = 2.0\)) и снижается в спортзале (\(\epsilon = 0.5\)). Температура \(\tau = 0.1\). Шанс отказа 20\%.
    
    \item \textbf{Нонконформист} — импульсивный и склонный к конфликту архетип. Часто действует наперекор ожиданиям, провоцируя эмоциональные всплески. 
    \textit{Особенности}: Формула приоритета \(P_{ij} = \epsilon \cdot (A_{ij} + U_{ij} + \sin(T_{ij}))\). Адаптивность к среде нейтральна (\(\epsilon = 1.0\)). Температура \(\tau = 1.5\). Шанс отказа 40\%.
    
    \item \textbf{Медиатор} — стремится к балансу и согласию. Склонен принимать предложения о взаимодействии, повышенно чувствителен к симпатии и доверию. 
    \textit{Особенности}: Формула приоритета \(P_{ij} = \epsilon \cdot (\text{sigmoid}(A_{ij}) + U_{ij} + \exp(T_{ij}))\). Адаптивность к среде максимальна на переменах (\(\epsilon = 2.0\)). Температура \(\tau = 0.3\). Шанс отказа всего 5\%.
    
    \item \textbf{Конкурент} — решительный и смелый, архетип с низкой чувствительностью к страху и отвращению. Быстро реагирует на угрозы, активно участвует в конфликтах. 
    \textit{Особенности}: Формула приоритета \(P_{ij} = \epsilon \cdot (A_{ij} + U_{ij} + \text{sigmoid}(T_{ij}))\). Адаптивность к среде максимальна в спортзале (\(\epsilon = 3.0\)). Температура \(\tau = 0.1\). Шанс отказа 20\%.
    
    \item \textbf{Оппортунист} — непредсказуем, легко переключается между эмоциями. Взаимодействует спонтанно, часто влияет на эмоциональные состояния других. 
    \textit{Особенности}: Формула приоритета \(P_{ij} = \epsilon \cdot (\sin(A_{ij}) + U_{ij} + T_{ij})\). Адаптивность к среде нейтральна (\(\epsilon = 1.0\)). Температура \(\tau = 3.0\). Шанс отказа 50\%.
    
    \item \textbf{Протектор} — охранительный архетип, заботится о других, склонен к устойчивым связям. Повышенно чувствителен к чувству долга. 
    \textit{Особенности}: Формула приоритета \(P_{ij} = \epsilon \cdot (A_{ij} + U_{ij} + \text{sigmoid}(T_{ij}))\). Адаптивность к среде нейтральна (\(\epsilon = 1.0\)). Температура \(\tau = 0.2\). Шанс отказа 10\%.
    
    \item \textbf{Деструктор} — сосредоточен на распаде, сомнении, исчезновении. Склонен разрушать устойчивые связи, вносит хаос в отношения. 
    \textit{Особенности}: Формула приоритета \(P_{ij} = \epsilon \cdot (\log(A_{ij}) + U_{ij} + \sin(T_{ij}))\). Адаптивность к среде нейтральна (\(\epsilon = 1.0\)). Температура \(\tau = 1.0\). Шанс отказа 80\%.
    
    \item \textbf{Инноватор} — уравновешен и открыт к новому, ищет смысл во взаимодействиях, инициирует связи на основе интеллектуальной и моральной оценки. 
    \textit{Особенности}: Формула приоритета \(P_{ij} = \epsilon \cdot (A_{ij} + \log(U_{ij}) + T_{ij})\). Адаптивность к среде нейтральна (\(\epsilon = 1.0\)). Температура \(\tau = 1.2\). Шанс отказа 20\%.
    
    \item \textbf{Рефлектор} — ориентирован на прошлое, устойчив к эмоциональным колебаниям, склонен к рефлексии и сохранению долгосрочных связей. 
    \textit{Особенности}: Формула приоритета \(P_{ij} = \epsilon \cdot (\exp(A_{ij}) + U_{ij} + T_{ij})\). Адаптивность к среде нейтральна (\(\epsilon = 1.0\)). Температура \(\tau = 1.0\). Шанс отказа 20\%.
\end{itemize}

\section{Формальная постановка задачи}

Рассмотрим множество агентов \( \mathcal{K} = \{\mathcal{A}_1, \mathcal{A}_2, \ldots, \mathcal{A}_n \} \). Каждому агенту присваивается вектор эмоциональных состояний:
\[
E_i = (e_{i1}, e_{i2}, \ldots, e_{i7}) \in \mathbb{E}_7^7,
\]
где \( \mathbb{E}_7 = [-30, 30] \subset \mathbb{Z} \) — допустимый целочисленный диапазон для каждой из осей эмоций.

Его отношение к каждому другому агенту \( \mathcal{A}_j \) задается набором:
\[
U_{ij}(t),\; A_{ij}(t),\; T_{ij}(t) \in \mathbb{E}_{R},
\]
где \( \mathbb{E}_{R} = [-100, 100] \subset \mathbb{Z} \).

На каждом шаге симуляции агент вычисляет приоритет взаимодействия со всеми доступными агентами. Функция приоритета является нелинейной суммой математических трансформаций (\(f_A, f_U, f_T\)), задаваемых архетипом:
\[
P_{ij} = \epsilon \cdot \Bigl( f_A(A_{ij}) + f_U(U_{ij}) + f_T(T_{ij}) \Bigr),
\]
где $\epsilon$ — контекст среды.

Вероятность выбора конкретного кандидата $j$ вычисляется с помощью функции Softmax, учитывающей параметр температуры $\tau$:
\[
\text{Prob}(j) = \frac{\exp(P_{ij} / \tau)}{\sum_k \exp(P_{ik} / \tau)}
\]

После выбора цели агент получает дискретный исход \( \sigma_{i,j}(t) \). Набор исходов \(\{\sigma\}\) представляет собой совокупность результатов коммуникации агента с его окружением на текущем такте. В случае отказа (\(\sigma_{i,j} = 0\)) инициатор и адресат получают негативные изменения реляционных функций в зависимости от своего архетипа.

Далее агент обновляет своё эмоциональное состояние согласно векторной функции:
\[
E_i(t+1) = f_i\Bigl(E_i(t),\; s_i,\, \bigl\{ U_{ij}(t),\, A_{ij}(t),\, T_{ij}(t),\, \sigma_{i,j}(t) \bigr\}_{j=1}^n \Bigr).
\]
А также обновляет значения реляционных функций:
\[
\begin{aligned}
U_{ij}(t+1) &= g^{(U)}_i\big(U_{ij}(t), s_i, E_i(t), E_j(t), \sigma_{i,j}(t)\big), \\
A_{ij}(t+1) &= g^{(A)}_i\big(A_{ij}(t), s_i, E_i(t), E_j(t), \sigma_{i,j}(t)\big), \\
T_{ij}(t+1) &= g^{(T)}_i\big(T_{ij}(t), s_i, E_i(t), E_j(t), \sigma_{i,j}(t)\big).
\end{aligned}
\]

Для моделирования внешних факторов и долгосрочных циклов введены дополнительные механизмы коррекции эмоционального фона:
\begin{itemize}
    \item \textbf{Воскресная ремиссия:} в конце каждого недельного цикла (воскресенье) все агенты подвергаются воздействию «семейного дня». Эмоциональный вектор каждого агента получает стохастическое смещение $\Delta E_i = (\delta, \delta, \ldots, \delta)$, где $\delta \in [-30, 30]$ — единое случайное целое число для всех семи осей. Данный механизм моделирует внешнее социальное влияние, не связанное с внутригрупповыми взаимодействиями.
    \item \textbf{Семестровая переинициализация:} в начале каждого академического семестра (1 сентября и 1 февраля) происходит полный сброс эмоциональных состояний коллектива. Вектор $E_i$ каждого агента переинициализируется случайными значениями из диапазона $[-30, 30]$. Это отражает адаптацию к новому учебному расписанию и смену контекста среды.
\end{itemize}

\section{Описание модели}

Основой взаимодействия в модели выступает принцип распределения агентов в сложной среде Университета. Университет включает в себя 5 факультетов бакалавриата и 6 направлений магистратуры. На бакалавриате обучаются 4 курса, а в магистратуре — 2 курса. Каждый курс бакалавриата разбит на 3 учебные группы (по 25 студентов в каждой), а магистратуры — на 2 группы (по 10 студентов в каждой). Академический цикл внутри семестра состоит из 6 учебных дней в неделю. Ежедневно предусмотрено 4 обязательные пары. Симуляция времени учитывает реалистичный календарный цикл с разделением на осенние и весенние семестры (длительностью по 5 месяцев каждый) и перерывами на каникулы. Глобальный временной горизонт одной полномасштабной симуляции составляет 8 семестров, что соответствует 4 академическим годам и покрывает полный цикл обучения бакалавриата.

В зависимости от своего личного расписания и склонности к пропускам (\(\chi_i\)), агенты физически распределяются по локациям: лекционные залы, малые семинарские комнаты, центральный коридор (холл), а также спортзал, если по расписанию предусмотрена плановая пара физической культуры.

Помимо 4 обязательных учебных слотов, введено понятие 5-го слота — добровольные спортивные секции (кружки), которые проходят в спортзале после окончания основных пар. Существует 3 независимые секции: бадминтон, футбол и баскетбол. Студенты принимают решение о посещении спортзала на основе своего параметра спортивности (\(\rho_i\)). Если порог пройден, агент направляется в спортзал и присоединяется к одной из секций. Выбор конкретной секции происходит стохастически с учетом равномерного распределения, что формирует три изолированных графа взаимодействия внутри единой локации.

Взаимодействие агентов базируется на процессе группообразования, который математически описывается парадоксом дней рождений. Формирование групп происходит следующим образом:
\begin{enumerate}
    \item Агент анализирует свое окружение. В аудиториях и в спортзале (внутри секций) агент рассматривает только тех агентов, кто находится непосредственно рядом с ним (топология соседей). В коридоре же агенты активно бродят, и их окружение случайно перемешивается на каждой перемене, позволяя самостоятельно находить новых собеседников.
    \item Агент начинает набирать группу для общения, рекурсивно выбирая соседей на основе вероятностей Softmax.
    \item Процесс расширения группы продолжается до тех пор, пока очередным выбором не станет агент, который \textit{уже} состоит в этой группе, либо пока алгоритм не выберет "STOP-маркер" (усредненный приоритет окружения).
\end{enumerate}

Внутри сформированной группы применяется строгая топология клики — каждый участник вступает во взаимодействие с каждым другим участником образованной группы, независимо от ее финального размера.

Инициализация связей учитывает социальную физику: студенты первого курса бакалавриата и первого курса магистратуры, не обучавшиеся здесь ранее, входят в среду без сформированных связей. Студенты старших курсов уже обладают сложившейся внутригрупповой и межгрупповой сетью контактов.

\section{Программная архитектура модели}

Одной из центральных задач дипломной работы являлась модернизация программного комплекса для обеспечения возможности моделирования сверхбольших коллективов.

\subsection{Гибридное вычислительное ядро}
Вычисления в многоагентных системах имеют квадратичную алгоритмическую сложность $O(N^2)$. Для устранения узких мест производительности было разработано высокопроизводительное вычислительное ядро на языке C++. Ядро использует многопоточность по стандарту OpenMP для параллельного расчета векторов состояний. 

Важнейшим шагом оптимизации стал архитектурный переход от чисел с плавающей точкой к дискретным целочисленным шкалам. Это критически снизило потребление оперативной памяти и ускорило передачу данных между Python и C++.

\subsection{Инфраструктура хранения Big Data}
Для глубокого ретроспективного анализа была внедрена колоночная СУБД ClickHouse. Реализована система асинхронного логирования, которая на каждом такте отправляет в базу данных полные снимки эмоциональных состояний и матриц взаимодействий. Дополнительно реализован реестр агентов, связывающий числовые идентификаторы ядра со смысловыми метками объектов.

\subsection{Фоновый режим и управление}
Управляющая логика была вынесена в независимый класс \texttt{SimulationSession}. Внедрение фонового режима (headless) и системы логирования метаданных дает возможность исследователю запускать параллельные пакетные тестирования сценариев.

\section{Описание экспериментов}

% TODO: Данный раздел оставлен пустым для будущего описания методологии и результатов экспериментов.

\section{Заключение}

% TODO: Заключение будет сформулировано после проведения вычислительных экспериментов.

\clearpage
\section{Список литературы}

\begin{enumerate}
    \item Кудрявцев В. Б., Алешин С. В., Подколзин А. С. Теория автоматов: учебник для бакалавриата и магистратуры. — 2-е изд., испр. и доп. — М.: Юрайт, 2019. — 320 с.
    \item Волков Н. Ю. Об автоматной модели преследования. — Дискретная математика, т.19, вып.2, 2007, стр. 131–160.
\end{enumerate}

\end{document}
